\chapter{The Smith-Waite Tarot}\label{the-smith-waite-tarot}

\section{Operating Principles}\label{operating-principles}

The Smith-Waite Tarot was designed by Pamela Colman Smith under the direction of Arthur Edward Waite, completed in 1909 in Smith\textquotesingle s Chelsea studio. Smith painted all 78 images in six months using ink and watercolor, working from Waite\textquotesingle s detailed instructions for the Major Arcana and from her own invention for the Minor Arcana pip cards. The pictorial scenes on the 56 Minor Arcana cards --- the single most significant development in tarot history --- were Smith\textquotesingle s creation, not Waite\textquotesingle s. Waite had little interest in the Minors. Smith transformed what had been geometric arrangements of suit symbols into a complete visual vocabulary for human experience.

Both were members of the Hermetic Order of the Golden Dawn. The esoteric framework of the deck is therefore the same Qabalistic structure underlying the Golden Dawn Magical Tarot --- Hebrew letters, Paths on the Tree of Life, elemental and astrological correspondences. Waite deliberately obscured some of this structure for non-initiates while retaining it intact beneath the surface. The deck was designed to function both as a divinatory tool for general use and as an esoteric instrument for those with the training to read its deeper layers. It succeeds at both simultaneously.

Smith\textquotesingle s monogram --- a serpentine snake entwined around her initials --- appears on every card. It is the only authorial mark on the deck and she placed it there herself.

This system uses reversals. An upright card and a reversed card carry different meanings. This distinguishes it from both the Golden Dawn system (which uses elemental dignities) and the English playing card system (which uses position and combination).

\begin{center}\rule{0.5\linewidth}{0.5pt}\end{center}

\section{The Structure}\label{the-structure}

78 cards in two groups:

\textbf{The Major Arcana} --- 22 keys numbered 0 through XXI. These follow the Golden Dawn Path structure: each key is a Path on the Tree of Life connecting two Sephiroth, carrying a Hebrew letter, an astrological attribution, and a specific force. Waite retained the GD framework but corrected one significant misalignment in the original Book T correspondence (see Departures section).

\textbf{The Minor Arcana} --- 56 cards in four suits of 14 cards each: Ace through 10, plus four court cards (King, Queen, Knight, Page). The 40 pip cards (Ace through 10) carry Smith\textquotesingle s pictorial scenes. The 16 court cards carry armored or robed figures whose expression, posture, and environment Smith used to convey character rather than merely rank.

\begin{center}\rule{0.5\linewidth}{0.5pt}\end{center}

\section{The Suits and Their Correspondences}\label{the-suits-and-their-correspondences}

\textbf{Wands --- Fire} Will, enterprise, ambition, the creative and spiritual impulse, energy before it finds form. The suit of what is being initiated or driven forward. Wands figures are generally depicted in landscapes of open sky and bare hills --- the unformed territory that fire moves through.

\textbf{Cups --- Water} Emotion, imagination, the interior life, relationship, the unconscious. The suit of feeling and vision. Cup scenes often feature water, calm or turbulent, and emphasize interiority --- figures contemplating, dreaming, receiving.

\textbf{Swords --- Air} Intellect, conflict, truth, the consequence of thought and word. The most difficult suit in the deck. Swords scenes frequently depict aftermath --- the storm has passed or is passing, figures standing in the wreckage of what clear thought or harsh truth has produced. Air both clarifies and cuts.

\textbf{Pentacles --- Earth} The material world, the body, money, craft, the practical result. The suit of what has actually manifested. Pentacle scenes tend toward ordered gardens, cultivated land, workshops --- the world where human effort meets physical reality.

\begin{center}\rule{0.5\linewidth}{0.5pt}\end{center}

\section{The Esoteric Framework}\label{the-esoteric-framework}

Waite retained the Golden Dawn Qabalistic structure beneath the surface of the deck. The 22 Major Arcana correspond to the 22 Paths on the Tree of Life. The 10 pip numbers per suit correspond to the 10 Sephiroth. The four suits correspond to the four worlds (Atziluth, Briah, Yetzirah, Assiah) and to the four letters of YHVH.

This framework is not printed on the cards and Waite deliberately withheld the correspondences from his written commentary to protect initiatory teaching from general publication. The Pictorial Key to the Tarot (1910), his companion volume, is notably evasive on these points while accurate on divinatory meanings. The full correspondence system must be drawn from the Golden Dawn documents directly.

For practical reference, the Sephirothic assignments of the pip cards are identical to those in the Golden Dawn system:

\begin{longtable}[]{@{}>{\RaggedRight\small}p{0.18\linewidth}>{\RaggedRight\small}p{0.22\linewidth}>{\RaggedRight\small}p{0.54\linewidth}@{}}
\toprule\noalign{}
Number & Sephirah & Character \\
\midrule\noalign{}
\endhead
\bottomrule\noalign{}
\endlastfoot
Ace & Kether & Root force; undivided, radical power of the suit \\
2 & Chokmah & Initial polarity; the first differentiation; beginnings \\
3 & Binah & First concrete form; realization through limitation \\
4 & Chesed & Stability; the foundation established \\
5 & Geburah & Disruption; the established order disturbed \\
6 & Tiphareth & Harmony; accomplished success; the solar center \\
7 & Netzach & Force seeking outlet; result dependent on action \\
8 & Hod & Swift directed force; precise but potentially hasty \\
9 & Yesod & Fundamental strength; the approach of completion \\
10 & Malkuth & Final manifestation; the matter fully determined \\
\end{longtable}

\begin{center}\rule{0.5\linewidth}{0.5pt}\end{center}

\section{Key Departures from the Golden Dawn Tradition}\label{key-departures-from-the-golden-dawn-tradition}

\textbf{The Strength/Justice correction.} In the original Golden Dawn Book T, the astrological correspondence assigned Leo (Teth) to Key VIII and Libra (Lamed) to Key XI. In the traditional Marseilles ordering, Key VIII is Justice and Key XI is Strength. This produced a mismatch: the Justice card (scales imagery) was assigned to Leo, and the Strength card (lion imagery) was assigned to Libra --- the reverse of what the iconography implies.

Waite corrected this by swapping the two cards:

\begin{itemize}
\tightlist
\item
  Key VIII = \textbf{Strength} = Teth = Leo (lion imagery now matches Leo attribution)
\item
  Key XI = \textbf{Justice} = Lamed = Libra (scales imagery now matches Libra attribution)
\end{itemize}

This swap is the most significant structural departure from both the Marseilles tradition and from the original GD numbering. The Golden Dawn Magical Tarot (Cicero) follows the same corrected order.

The Marseilles tradition: Justice = VIII, Strength = XI. Smith-Waite and Cicero GD: Strength = VIII, Justice = XI.

When working across decks, this is the single most important structural difference to track.

\textbf{The pictorial Minor Arcana.} All tarot decks before Smith-Waite used geometric arrangements of suit symbols for the pip cards --- four cups arranged in a pattern for the Four of Cups, eight swords for the Eight of Swords. There were no narrative scenes. Smith introduced pictorial scenes to all 40 pip cards (Ace through Ten in each suit). These scenes were her invention, developed from Waite\textquotesingle s suit meanings with considerable creative latitude. The Minor Arcana as a pictorial narrative vocabulary is Smith\textquotesingle s contribution to the tradition, not Waite\textquotesingle s.

The consequence for reading practice is significant: Smith\textquotesingle s scenes communicate meaning directly through image, bypassing the need to know the underlying Qabalistic or astrological structure. A reader who knows nothing of Sephiroth can read the Five of Swords (three figures, two retreating in defeat, one gathering the swords with a cold expression) and derive its meaning accurately from the scene alone. This is what made the deck accessible to non-initiates and what drove its dominance throughout the twentieth century.

\textbf{Court card naming.} Waite reverted to traditional pre-GD naming for the court cards: King, Queen, Knight, Page. The GD used King, Queen, Prince, Princess to reflect the YHVH tetrad (Yod, Heh, Vau, Heh final). Waite\textquotesingle s naming obscures the elemental-rank correspondence:

\begin{itemize}
\tightlist
\item
  Smith-Waite \textbf{Knight} (mounted) corresponds to GD \textbf{King} (Yod, Fire of the suit)
\item
  Smith-Waite \textbf{King} (seated) corresponds to GD \textbf{Prince} (Vau, Air of the suit)
\item
  Smith-Waite \textbf{Queen} = GD \textbf{Queen} (Heh, Water of the suit)
\item
  Smith-Waite \textbf{Page} = GD \textbf{Princess} (Heh final, Earth of the suit)
\end{itemize}

This means the most senior court card by elemental rank in the GD system --- the Yod, Fire quality --- is called the Knight in Smith-Waite, while the card called King is actually the third in elemental seniority. Waite did not document this alignment clearly and modern readers largely ignore it, reading the courts by scene and character rather than elemental rank.

\textbf{Renaming of Major Arcana.} Waite renamed several cards from their traditional titles: the Papess became the High Priestess, the Pope became the Hierophant. These renamings removed Catholic imagery in favor of broader esoteric terminology consistent with the GD\textquotesingle s non-denominational magical framework.

\begin{center}\rule{0.5\linewidth}{0.5pt}\end{center}

\section{The Major Arcana}\label{the-major-arcana}

\begin{longtable}[]{@{}>{\RaggedRight\small}p{0.08\linewidth}>{\RaggedRight\small}p{0.12\linewidth}>{\RaggedRight\small}p{0.10\linewidth}>{\RaggedRight\small}p{0.10\linewidth}>{\RaggedRight\small}p{0.12\linewidth}>{\RaggedRight\small}p{0.42\linewidth}@{}}
\toprule\noalign{}
Key & Name & Hebrew Letter & Attribution & Path & Waite\textquotesingle s Title \\
\midrule\noalign{}
\endhead
\bottomrule\noalign{}
\endlastfoot
0 & The Fool & Aleph & Air & 11 (Kether--Chokmah) & Spirit of the Ether \\
I & The Magician & Beth & Mercury & 12 (Kether--Binah) & Magus of Power \\
II & The High Priestess & Gimel & Moon & 13 (Kether--Tiphareth) & Priestess of the Silver Star \\
III & The Empress & Daleth & Venus & 14 (Chokmah--Binah) & Daughter of the Mighty Ones \\
IV & The Emperor & Heh & Aries & 15 (Chokmah--Tiphareth) & Son of the Morning \\
V & The Hierophant & Vau & Taurus & 16 (Chokmah--Chesed) & Magus of the Eternal \\
VI & The Lovers & Zayin & Gemini & 17 (Binah--Tiphareth) & Children of the Voice \\
VII & The Chariot & Cheth & Cancer & 18 (Binah--Geburah) & Child of the Powers of the Waters \\
VIII & Strength & Teth & Leo & 19 (Chesed--Geburah) & Daughter of the Flaming Sword \\
IX & The Hermit & Yod & Virgo & 20 (Chesed--Tiphareth) & Prophet of the Eternal \\
X & Wheel of Fortune & Kaph & Jupiter & 21 (Chesed--Netzach) & Lord of the Forces of Life \\
XI & Justice & Lamed & Libra & 22 (Geburah--Tiphareth) & Daughter of the Lords of Truth \\
XII & The Hanged Man & Mem & Water & 23 (Geburah--Hod) & Spirit of the Mighty Waters \\
XIII & Death & Nun & Scorpio & 24 (Tiphareth--Netzach) & Child of the Great Transformers \\
XIV & Temperance & Samekh & Sagittarius & 25 (Tiphareth--Yesod) & Daughter of the Reconcilers \\
XV & The Devil & Ayin & Capricorn & 26 (Tiphareth--Hod) & Lord of the Gates of Matter \\
XVI & The Tower & Peh & Mars & 27 (Netzach--Hod) & Lord of the Hosts of the Mighty \\
XVII & The Star & Tzaddi & Aquarius & 28 (Netzach--Yesod) & Daughter of the Firmament \\
XVIII & The Moon & Qoph & Pisces & 29 (Netzach--Malkuth) & Ruler of Flux and Reflux \\
XIX & The Sun & Resh & Sol & 30 (Hod--Yesod) & Lord of the Fire of the World \\
XX & Judgement & Shin & Fire & 31 (Hod--Malkuth) & Spirit of the Primal Fire \\
XXI & The World & Tau & Saturn & 32 (Yesod--Malkuth) & Great One of the Night of Time \\
\end{longtable}

\subsection{Major Arcana: Smith\textquotesingle s Scenes and Waite\textquotesingle s Meanings}\label{major-arcana-smiths-scenes-and-waites-meanings}

\textbf{0 --- The Fool} \emph{Scene:} A young figure in ornate clothing steps toward the edge of a cliff, eyes upward, carrying a white rose and a small pack on a wand. A small dog leaps at his heels. Mountains behind. Sky clear. \emph{Upright:} Folly, mania, extravagance in the best sense --- the leap into the unknown that is also pure spiritual potential. The beginning before beginning. Unconditioned. \emph{Reversed:} Negligence, apathy, the leap made without attention; foolishness without redemption.

\textbf{I --- The Magician} \emph{Scene:} A figure in red and white robes stands before a table bearing the four suit symbols. One hand raises a wand toward heaven; one points to earth. Above his head, the lemniscate (infinity sign). About his waist, a serpent devouring its own tail. \emph{Upright:} Skill, will, self-confidence, the capacity to act. The directed use of all available resources. Initiation of a new cycle. \emph{Reversed:} Will misdirected; skill used for manipulation; the appearance of power without its substance.

\textbf{II --- The High Priestess} \emph{Scene:} A crowned figure sits between two pillars --- one black (Boaz), one white (Jachin) --- before a veil hung with pomegranates. A crescent moon at her feet. She holds a scroll marked TORA, partially concealed. \emph{Upright:} Hidden knowledge, the unconscious, what is not yet revealed, inner wisdom accessible only through stillness. The unrealized future. \emph{Reversed:} Superficial knowledge; what is hidden used for manipulation; passion overriding wisdom.

\textbf{III --- The Empress} \emph{Scene:} A crowned, pregnant figure sits in a field of ripe grain before a forest. Heart-shaped shield bearing the Venus symbol. Flowing robes covered in pomegranates. \emph{Upright:} Fertility, abundance, nature, the creative force in full productive expression. Growth, sensory pleasure, domestic life flourishing. \emph{Reversed:} Inaction, dissipation, creative energy without outlet or squandered.

\textbf{IV --- The Emperor} \emph{Scene:} An armored, bearded figure sits rigidly on a stone throne decorated with rams\textquotesingle{} heads. He holds an ankh and an orb. Mountains behind --- hard and red. \emph{Upright:} Temporal power, governance, the establishing of order through will and force. Authority, the father. Dominion over the material world. \emph{Reversed:} Domination, rigidity, the abuse of authority; weakness concealed behind show of force.

\textbf{V --- The Hierophant} \emph{Scene:} A robed, crowned figure sits between two pillars, right hand raised in blessing. Two supplicants kneel before him. Two crossed keys at his feet. \emph{Upright:} Established institutions and their wisdom, teaching, the transmission of sacred knowledge, religious or formal initiation. The outer form of inner truth. \emph{Reversed:} Dogma mistaken for wisdom; conformity; the institution protecting itself rather than serving.

\textbf{VI --- The Lovers} \emph{Scene:} A man and woman stand naked before an angel in the sky --- Adam and Eve before the angel Raphael. Behind the woman, the Tree of Knowledge. Behind the man, a tree of flame. Mountain in the distance. \emph{Upright:} The right choice made through love and aspiration rather than calculation. Union, harmony, alignment of values. The beginning of genuine partnership. \emph{Reversed:} Wrong choice; separation; values in conflict with each other.

\textbf{VII --- The Chariot} \emph{Scene:} A crowned warrior stands in a stone chariot drawn by two sphinxes --- one black, one white --- in apparent opposition. He carries a wand and is armored. Stars above, city behind. \emph{Upright:} Triumph through the mastery of opposing forces. Victory, discipline, the will directing contradictory energies toward a single end. Conquest. \emph{Reversed:} Defeat by one\textquotesingle s own contradictions; victory sought without ethical grounding.

\textbf{VIII --- Strength} \emph{Scene:} A woman in white, crowned with flowers, gently closes the mouth of a lion. The lemniscate above her head. No struggle --- she holds the beast with ease. \emph{Upright:} Courage, patience, the gentle mastery of one\textquotesingle s own animal nature. Force that does not need violence. Endurance. The strength that comes from inner certainty. \emph{Reversed:} Weakness, abuse of power, desire overriding judgment.

\textbf{IX --- The Hermit} \emph{Scene:} An old, robed figure stands alone on a snowy peak. He holds a lantern containing a six-pointed star --- the Seal of Solomon. He carries a staff. \emph{Upright:} Wisdom sought and found in solitude. The guide who has completed the journey and now carries the light for others. Prudence. Withdrawal in service of illumination. \emph{Reversed:} Excessive withdrawal, fear disguised as wisdom; refusing the return.

\textbf{X --- Wheel of Fortune} \emph{Scene:} A large wheel inscribed with alchemical and Qabalistic symbols turns in space. A sphinx sits on top; below, Typhon descends. Four winged figures at the corners hold books --- the symbols of the four fixed signs (man, eagle, lion, bull). \emph{Upright:} Fortune turning, good luck arriving, the completion of a cycle. What rises when the wheel moves. The larger pattern asserting itself. \emph{Reversed:} Fortune turning downward; bad luck; resistance to necessary change.

\textbf{XI --- Justice} \emph{Scene:} A crowned, robed figure sits between two pillars, holding scales in one hand and an upraised sword in the other. The figure is still, formal, impartial. \emph{Upright:} Equity, righteous judgment, the truth of the situation as it actually stands. What is fair. Legal matters decided correctly. The equilibrating principle. \emph{Reversed:} Bias, legal injustice, judgment corrupted by interest.

\textbf{XII --- The Hanged Man} \emph{Scene:} A young man hangs by one foot from a T-shaped wooden frame between two living trees. His expression is calm, almost radiant. A halo of light surrounds his head. His free leg crosses behind the suspended leg, forming a figure four. \emph{Upright:} Willing sacrifice, suspension, the reversal of ordinary values that allows a higher perspective. What cannot be forced. The necessary pause. \emph{Reversed:} Arrogance in sacrifice; the martyrdom that serves the ego; wasted sacrifice.

\textbf{XIII --- Death} \emph{Scene:} An armored skeleton on a white horse moves through a field. A king lies dead before it. A bishop approaches in prayer. A young woman turns away. A child offers flowers. In the distance, the sun rises between two towers. \emph{Upright:} Transformation, the necessary end. What must die for what is next to begin. The card rarely means physical death --- it means the death of a situation, identity, or phase of life. \emph{Reversed:} Stagnation, the refusal to let something die that has finished; inertia.

\textbf{XIV --- Temperance} \emph{Scene:} An angel stands with one foot on land, one in water, pouring liquid between two cups in a continuous stream. Irises grow near the water. A path leads toward mountains and a shining crown in the sky. \emph{Upright:} The combination of forces into a working whole. Moderation as an active art. The alchemical mixing that produces something neither ingredient contained. Patience, adaptation, balance in motion. \emph{Reversed:} Discord, competing forces that will not combine; excess.

\textbf{XV --- The Devil} \emph{Scene:} A goat-headed figure sits on a black cube to which two human figures are chained by loose chains --- they could remove them. The figures have horns and tails. The Devil holds an inverted torch. \emph{Upright:} Bondage to the material plane, to habit, to compulsion. The chains are chosen, not locked. The force of matter fully realized --- which may serve or bind depending on how it is engaged. \emph{Reversed:} Beginning of release from bondage; the first recognition of the chain.

\textbf{XVI --- The Tower} \emph{Scene:} Lightning strikes a tall stone tower. A crown is blasted from its top. Two figures fall headlong. Flames issue from the windows. \emph{Upright:} Sudden catastrophic change. The structure built on false foundations is struck down. Cannot be avoided once in motion. Sometimes liberation through destruction. \emph{Reversed:} The catastrophe suppressed, postponed, built up further; oppression.

\textbf{XVII --- The Star} \emph{Scene:} A naked woman kneels at the edge of water, pouring from two vessels --- one into the water, one onto the land. Eight stars above her, one large and central. \emph{Upright:} Hope, faith, spiritual nourishment freely given. The grace that arrives after the darkness of The Tower. Clarity, inspiration, unexpected help. \emph{Reversed:} Hope deferred; self-doubt masquerading as realism; gifts withheld.

\textbf{XVIII --- The Moon} \emph{Scene:} A full moon with a face looks down on two towers. A wolf and a dog howl at it from either side of a path that leads into the distance. A crayfish emerges from a pool. \emph{Upright:} Illusion, deception, the terror and fertility of the unconscious. The liminal state between worlds. Things are not as they appear. The path must be walked without full sight. \emph{Reversed:} Deception recognized too late; practical mistakes; confusion less than feared.

\textbf{XIX --- The Sun} \emph{Scene:} A large radiant sun shines down on a naked child riding a white horse before a wall hung with sunflowers. The child holds a red banner. \emph{Upright:} Success, contentment, clarity, the pleasure of being fully alive in a good world. The most unambiguously fortunate of the Major Arcana. \emph{Reversed:} Success delayed or diminished; happiness present but obscured.

\textbf{XX --- Judgement} \emph{Scene:} An angel blows a trumpet from the clouds. Figures rise from coffins below --- men, women, children --- with arms raised. A mountain range in the distance. \emph{Upright:} The call that cannot be ignored. The final reckoning with one\textquotesingle s own past. A decisive outcome. The awakening to what one has actually been. \emph{Reversed:} The call heard but refused; delay in reckoning; the awakening postponed.

\textbf{XXI --- The World} \emph{Scene:} A naked figure dances within a laurel wreath, holding two wands. Four winged figures at the corners --- bull, lion, eagle, angel --- the fixed signs. The wreath is oval, enclosing the dancer completely. \emph{Upright:} Synthesis, completion, the successful conclusion of a long cycle. The world as it actually is --- whole. Often describes the situation itself rather than adding a force to it. \emph{Reversed:} Incomplete success; the final step not yet taken; stagnation at the threshold of completion.

\begin{center}\rule{0.5\linewidth}{0.5pt}\end{center}

\section{The Aces}\label{the-aces}

The Aces are Kether of their suit --- the root force before it differentiates into the ten Sephiroth below. In Smith\textquotesingle s images, each Ace presents a single hand emerging from a cloud, offering the suit symbol. The hand is the hand of the divine --- impersonal, neither male nor female, offering the instrument to whoever will take it.

\textbf{Ace of Wands} A hand offers a single flowering wand. Flames leave its tip. A castle in the distance. \emph{Upright:} The beginning of a creative or spiritual enterprise. Raw initiative. The spark that precedes the fire. Fortune and enterprise. \emph{Reversed:} Delays in starting; false start; the initiative that collapses before it builds.

\textbf{Ace of Cups} A chalice overflows into a pool below. A dove descends with a communion wafer. The cup has five streams pouring from it. \emph{Upright:} The beginning of love, abundance, joy, spiritual nourishment freely offered. The overflowing source. A new emotional beginning. \emph{Reversed:} Love withdrawn; the cup offered but refused; spiritual emptiness.

\textbf{Ace of Swords} A hand holds a sword upright, crowned at its tip. Mountains below, barren. The crown is circled with the wreath of victory, trailing olive and palm. \emph{Upright:} The invoked power of the mind --- clarity, triumph, the force of truth. When raised toward heaven, the sword calls down divine clarity. The double edge: the same force liberates or destroys depending on how it is held. \emph{Reversed:} The blade turned against its bearer. Tyranny, disaster, the force of intellect misused.

\textbf{Ace of Pentacles} A hand offers a single large pentacle against a garden of white lilies and roses. A gate leads through a hedge into a mountainous distance. \emph{Upright:} Material prosperity, the beginning of a practical achievement, the seed of physical manifestation. The gate to the material garden is open. \emph{Reversed:} Prosperity blocked or misused; material opportunity refused.

\begin{center}\rule{0.5\linewidth}{0.5pt}\end{center}

\section{The Court Cards}\label{the-court-cards}

Court cards represent people in the querent\textquotesingle s life, the querent in a particular mode, or dominant forces shaping the situation. Smith gave each court figure a specific environment, expression, and posture that communicates character directly. The elemental assignments follow the same logic as the Golden Dawn system --- see the Departures section for the naming correspondence.

\subsection{Kings --- Seated figures; established authority; the mature expression of the suit\textquotesingle s force}\label{kings--seated-figures-established-authority-the-mature-expression-of-the-suits-force}

\textbf{King of Wands} Enthroned in the open air, armored, holding a flowering wand. Salamanders on his robe. A lion on the throne. He leans forward slightly --- alert, not resting. A generous, fiery, noble person. Quick in action and affection, impulsive, passionate. Natural authority that burns fast and moves on. \emph{Reversed:} Severity, prejudice, or the fire turned to arrogance.

\textbf{King of Cups} Enthroned on a stone block in the middle of a turbulent sea. He holds a cup and a short scepter. A fish leaps from the sea behind him. He appears calm amid the storm. A man of feeling who has mastered his interior life --- creative, skilled, emotionally intelligent. Commands without suppression. \emph{Reversed:} Craftiness and unscrupulousness; emotional manipulation.

\textbf{King of Swords} Enthroned, facing forward, armored, sword raised. His throne is decorated with butterflies and sylphs. The sky behind him is overcast with wind-bent trees. A figure of keen intellect and authority. Analytical, decisive, active in matters requiring judgment. Can be severe, even ruthless when convinced of being right. \emph{Reversed:} Cruelty, perversity, deliberate harm.

\textbf{King of Pentacles} Enthroned in a garden. His robe is covered with grapes and flowers. His throne is decorated with bulls\textquotesingle{} heads. He holds a pentacle and a scepter. Reliable, established, successful in material domains. Patient, methodical, not easily moved. His wealth is earned and maintained rather than won. The steward. \emph{Reversed:} Vice, weakness, corruption of what was solid.

\subsection{Queens --- Seated figures; receptive, interior power; the suit\textquotesingle s force made stable}\label{queens--seated-figures-receptive-interior-power-the-suits-force-made-stable}

\textbf{Queen of Wands} Enthroned before sunflowers. She holds a sunflower in one hand, a wand in the other. A black cat sits at her feet. She faces outward --- open, confident. Practical and fiery. Self-directed, passionate, attractive in personality. Loyal when her loyalty is earned. The most independent of the queens. \emph{Reversed:} Jealousy, fickleness, opposition wearing a friendly face.

\textbf{Queen of Cups} Enthroned at the water\textquotesingle s edge, holding an elaborate closed cup that she contemplates. Her robe flows into the sea. She does not look at the observer. Profoundly intuitive, visionary, emotionally deep. A person who receives impressions from the unconscious and can communicate them. The interior life is her home. \emph{Reversed:} The same depths become perversity; the vision turns to delusion.

\textbf{Queen of Swords} Enthroned on a cloud-height throne. Sword raised, left hand extended as if beckoning. Her expression is stern but not unkind. Sylphs in the clouds around her. A person who has experienced loss and sees clearly as a result. Keen, just, quick in perception. May have suffered, but the suffering has produced precision rather than bitterness. \emph{Reversed:} Malice, narrow-mindedness; sorrow used as a weapon.

\textbf{Queen of Pentacles} Enthroned in a garden surrounded by flowers and fruit. A rabbit runs at the lower corner. She holds a pentacle and contemplates it with gentleness. Practical, nurturing, resourceful. The ability to make any environment fertile. Generous with what she has. Entirely at home in the physical world. \emph{Reversed:} Neglect of the domestic world; suspicion; materialism without generosity.

\subsection{Knights --- Mounted figures; active, initiating force; the suit in motion}\label{knights--mounted-figures-active-initiating-force-the-suit-in-motion}

\textbf{Knight of Wands} On a rearing horse in an arid landscape, armored in red and gold with salamanders on his tunic. He charges forward, wand raised. Energetic, impulsive, departure from what is familiar, travel, a sudden journey. The fire force in full movement --- brilliant and unreliable by turns. \emph{Reversed:} Interruption, discord, the charge that leads nowhere.

\textbf{Knight of Cups} On a slow-walking horse in a calm landscape, armored, holding a cup forward as if presenting it. Wings on his helmet. River and fish in the distance. A romantic approach, an invitation, a proposition. Someone who brings feeling forward rather than acting immediately. The arrival of a love interest or a creative offer. \emph{Reversed:} Fraud, trickery; the romantic gesture that conceals manipulation.

\textbf{Knight of Swords} On a charging horse, armor flying, sword raised. The sky behind him is turbulent, trees bent by wind. He moves with absolute speed and commitment. Forceful, brave, domineering, impulsive in the extreme. The mind in full charge. Can be heroic or catastrophically ill-timed depending on circumstances. \emph{Reversed:} Imprudence, extravagance, the charge that destroys everything including the purpose it was meant to serve.

\textbf{Knight of Pentacles} On a stationary heavy horse in a plowed field. He holds a pentacle and studies it. Armored in dull greens and browns. Nothing is moving. Responsible, patient, methodical, dependable to the point of immobility. The person who will actually complete the task correctly and without drama. \emph{Reversed:} Stagnation, idleness, materialism as an end in itself.

\subsection{Pages --- Standing figures; nascent energy; new beginnings and messages}\label{pages--standing-figures-nascent-energy-new-beginnings-and-messages}

\textbf{Page of Wands} Stands in an arid landscape holding a long wand, looking at it with curiosity. His cap has a feather. Young, attentive, fired with new possibility. A bearer of news, a faithful young person, a creative impulse not yet in motion. Something is about to begin; the messenger has arrived. \emph{Reversed:} Bad news; an unstable young person; a project that will not begin.

\textbf{Page of Cups} Stands at the water\textquotesingle s edge holding a cup from which a fish has emerged and appears to be speaking to him. He regards it with gentle surprise. Reflective news, a message with emotional content. A gentle, dreamy young person or the beginning of an imaginative project. The unexpected that arrives from the interior. \emph{Reversed:} Deception, seduction; the dream that misleads.

\textbf{Page of Swords} Stands on a windswept cliff, sword raised, looking alertly about. His posture is ready, slightly suspicious. Vigilance, insight, the need to be alert. A perceptive young person. News involving scrutiny or negotiation. The arrival of a situation requiring sharp attention. \emph{Reversed:} The same vigilance turned to cunning or spying.

\textbf{Page of Pentacles} Stands in a garden holding a pentacle up before him and gazing at it with concentration. Flowers at his feet. Young, serious, absorbed. Practical news, the beginning of a material project, a diligent young person with capacity for growth. Scholarship or apprenticeship. \emph{Reversed:} Dissipation, wastefulness, the practical gift squandered.

\begin{center}\rule{0.5\linewidth}{0.5pt}\end{center}

\section{The Minor Arcana: Pip Cards 2--10}\label{the-minor-arcana-pip-cards-210}

Smith\textquotesingle s scenes are not arbitrary decoration --- they are pictorial statements of the card\textquotesingle s force, derived from the Sephirothic and astrological attributions underlying each card. Reading the scene is reading the force. The upright meaning follows the scene\textquotesingle s narrative; the reversed meaning is generally the same force turned inward, blocked, or distorted.

\subsection{Wands}\label{wands}

\begin{longtable}[]{@{}>{\RaggedRight\small}p{0.14\linewidth}>{\RaggedRight\small}p{0.18\linewidth}>{\RaggedRight\small}p{0.20\linewidth}>{\RaggedRight\small}p{0.42\linewidth}@{}}
\toprule\noalign{}
Card & Smith\textquotesingle s Scene & Upright & Reversed \\
\midrule\noalign{}
\endhead
\bottomrule\noalign{}
\endlastfoot
2 & A figure on a castle wall holds a globe, gazing at the sea. A second wand is fixed beside him. & Will established, planning the next step from a position of strength. Personal power, boldness. & Surprise, fear of the unknown, anticipation turning to anxiety. \\
3 & A figure with his back to us watches ships depart from a high place. Three wands are planted. & Enterprise under way. Practical success. The work in motion, the ships launched. & Caution, disappointment, beware of pride. \\
4 & Garlands hang between four wands. Two figures approach in celebration. A castle behind. & Completion, harvest, the rest that follows genuine effort. Prosperity, harmony, a settled situation. & Prosperity and happiness not fully claimed; beauty and completion withheld slightly. \\
5 & Five figures in open conflict, each wielding a wand against the others. & Strife, competition, conflict of interests. Not catastrophe --- the friction of opposed wills. & Contradiction, litigation, disputes. \\
6 & A mounted figure in laurel crown moves through a crowd, wand held high. Others hold wands alongside. & Victory, public recognition, good news. The successful outcome of a contest. & Apprehension of failure, disloyalty, the crown not yet secure. \\
7 & A figure defends a high position against six wands rising from below. & Valor, the courage to hold a position against opposition. Advantage from strength. & Perplexity, anxiety, embarrassment, being overwhelmed. \\
8 & Eight wands fly through open air over a landscape --- in transit, pointing toward a landing. & Movement, swift action, things arriving. News in transit. Haste. & Arrows of jealousy, quarrels, delays. \\
9 & A battered figure leans on a wand, guarding a row of eight behind him. He looks wary. & Great strength in reserve. The capacity to endure. Preparedness for opposition. & Obstacles, adversity, difficulty. \\
10 & A stooped figure carries ten wands, unable to see past them. & Oppression, the burden of too much taken on. Strength misapplied, force working against itself. & Difficulties, intrigues, the burden that is about to be set down. \\
\end{longtable}

\subsection{Cups}\label{cups}

\begin{longtable}[]{@{}>{\RaggedRight\small}p{0.14\linewidth}>{\RaggedRight\small}p{0.18\linewidth}>{\RaggedRight\small}p{0.20\linewidth}>{\RaggedRight\small}p{0.42\linewidth}@{}}
\toprule\noalign{}
Card & Smith\textquotesingle s Scene & Upright & Reversed \\
\midrule\noalign{}
\endhead
\bottomrule\noalign{}
\endlastfoot
2 & Two figures face each other, each holding a cup. A caduceus with a lion\textquotesingle s head rises between them. & Love, partnership, harmony between two people. The beginning of a true union. & Passion, unsatisfied desire, the separation of lovers. \\
3 & Three women dance and raise cups in celebration amid abundance. & Abundance, healing, joy shared among friends. Solace, pleasure, success in its social dimension. & Overindulgence, excess, pleasures beyond their appropriate limit. \\
4 & A figure sits beneath a tree, arms crossed, looking at three cups before them. A fourth cup is offered by a hand from a cloud --- unnoticed. & Weariness, discontent, satiation. The gift being offered that is not seen. Re-evaluation. & New possibilities, new relationships, new instructions. \\
5 & A cloaked figure stands before three spilled cups. Two upright cups stand behind, ignored. & Loss in what was pleasured. Regret, disillusionment. But not total loss --- two cups remain. & News, returning alliances, favorable expectations. \\
6 & A young figure gives flowers in a cup to a smaller figure in an old garden. & Looking to the past; nostalgia, memories, happiness from childhood. What was given freely. & The past as obstacle; living in memory; what is left behind. \\
7 & A figure gazes at seven cups floating in cloud, each containing a different vision --- treasure, laurels, a castle, a dragon, a figure of light, a wreath, a skull. & Fantasy, illusory success, the confusion of attractive visions. Many possibilities without the capacity to choose. & Determination, will applied to select among options. \\
8 & A figure walks away through the night from eight stacked cups, moving toward distant mountains. A moon above. & Abandoned success. The deliberate turning away from what has been built when it no longer satisfies. & Great joy, feasting, the return to what was left. \\
9 & A well-dressed figure sits with arms crossed before nine cups displayed on a shelf behind them. & Material happiness, the wish fulfilled. Contentment, physical well-being, satisfaction. & Imperfections, mistakes in the apparent success; truth withheld. \\
10 & Two adults stand embracing, arms raised toward a rainbow of ten cups. Two children dance nearby. & Perfected success in happiness. Family, home, lasting contentment --- the cycle completed. & Repose disturbed; family quarrels; the contentment incomplete. \\
\end{longtable}

\subsection{Swords}\label{swords}

\begin{longtable}[]{@{}>{\RaggedRight\small}p{0.14\linewidth}>{\RaggedRight\small}p{0.18\linewidth}>{\RaggedRight\small}p{0.20\linewidth}>{\RaggedRight\small}p{0.42\linewidth}@{}}
\toprule\noalign{}
Card & Smith\textquotesingle s Scene & Upright & Reversed \\
\midrule\noalign{}
\endhead
\bottomrule\noalign{}
\endlastfoot
2 & A blindfolded figure sits with arms crossed, holding two swords across their chest. A calm sea behind. & Stalemate, tension held in balance by will rather than resolution. Indecision, truce. & Release from the stalemate, movement in either direction, disloyalty revealed. \\
3 & Three swords pierce a heart suspended in a rainy sky. & Sorrow, removal, absence, delay, division. The mind registering a genuine wound. & Mental alienation, confusion, grief beginning to pass. \\
4 & A carved figure lies on a tomb in prayer position. Three swords hang on the wall; one is horizontal beneath the figure. & Repose from strife, rest, recuperation, solitude chosen for healing. & Renewed activity, occupation, caution in handling situations. \\
5 & A figure gathers three swords with a cold expression. Two others walk away in defeat. & Defeat, failure, hollow victory, the battle won at a cost that makes the winning meaningless. & Burial, obsequies, mourning, the aftermath of loss. \\
6 & A ferryman poles a boat across calm water. A figure with a child hunches in the bow, backs turned. Six swords stand upright in the bow. & Passage toward calmer waters. Earned movement through difficulty. Travel, gradual improvement. & Declaration, confession, the journey undertaken in difficulty. \\
7 & A figure carries five swords away from a camp, glancing back. Two swords remain planted. & Cunning, evasion, an attempt to succeed by stealth. Unstable effort, partial success through dishonesty. & Good advice, instruction, slander, the plan that fails. \\
8 & A bound and blindfolded figure stands surrounded by eight swords planted in a circle. A castle in the distance. & Crisis, restriction, interference --- imposed from without but maintained from within. Constraint. & Difficulty, treachery, the restriction beginning to loosen. \\
9 & A figure sits up in bed, face in hands, nine swords hanging on the wall behind in the darkness. & Desolation, despair, cruelty --- suffered or wielded. The darkest hour of the interior. & Imprisonment, suspicion, the nightmare real or near. \\
10 & A figure lies face down, ten swords in their back. Night behind; a thin light on the horizon. & Complete ending, the worst having happened. Also: the light on the horizon. Ruin, pain, but also the fact that it is done. & Advantage, profit, power --- arising from the ruins; what comes after. \\
\end{longtable}

\subsection{Pentacles}\label{pentacles}

\begin{longtable}[]{@{}>{\RaggedRight\small}p{0.14\linewidth}>{\RaggedRight\small}p{0.18\linewidth}>{\RaggedRight\small}p{0.20\linewidth}>{\RaggedRight\small}p{0.42\linewidth}@{}}
\toprule\noalign{}
Card & Smith\textquotesingle s Scene & Upright & Reversed \\
\midrule\noalign{}
\endhead
\bottomrule\noalign{}
\endlastfoot
2 & A figure balances two pentacles in a figure-eight loop of ribbon, while two ships ride rough seas behind. & Harmonious change, the juggling of two pursuits. Adaptability in difficult circumstances. & Enforced gaiety, simulated enjoyment, the balance failing. \\
3 & A sculptor works on a carving in a cathedral archway. Two figures consult a document before him. & Material works, craftsmanship, trade. The work being done with skill and recognized. & Mediocrity, ease without effort, the work below its potential. \\
4 & A crowned figure sits holding a pentacle to their chest, another pentacle below each foot, one on their head. The city behind. & Consolidated material power. Holding what has been gained. Security through careful management. Also: miserliness. & Delay, obstruction to material progress, opposition. \\
5 & Two ragged figures pass through snow beneath a lit church window. & Material trouble, poverty, the loss of what sustains. The spiritual resource goes untapped. & Chaos, disorder, money troubles --- but the worst may be passing. \\
6 & A figure with scales distributes coins to two kneeling supplicants. & Material success, generosity, the right flow of resources. Gifts given or received. & Desire, cupidity, debt, the generosity conditional. \\
7 & A figure leans on a hoe, looking at seven pentacles growing on a vine. & Success not yet gained, the wait for what has been planted. Ingenuity, labor, but the result still pending. & Cause for anxiety regarding money; impatience. \\
8 & A craftsman sits alone at a workbench, carefully making pentacles one by one. His work hangs behind him. & Apprenticeship, skilled labor, the mastery of a craft through patient application. Prudence. & Voided ambition, vanity, illegitimate work. \\
9 & A well-dressed figure stands in a garden among vines laden with grapes and pentacles. A hooded falcon on their wrist. & Accomplished material success through one\textquotesingle s own effort. Prudence, safety, well-being. & Roguery, dangerous acquaintance, the comfort built on uncertain foundations. \\
10 & An old man sits before a gate where two dogs rest. A family passes through the arch beyond. Ten pentacles in an arrangement above. & Wealth, the accumulation of what a life of effort builds. Legacy, family prosperity, the house in order. & Chance, fatality, the family fortune at risk. \\
\end{longtable}

\begin{center}\rule{0.5\linewidth}{0.5pt}\end{center}

\section{On Reversals}\label{on-reversals}

Reversals are used in this system. An upright card expresses its principal meaning. A reversed card generally indicates:

\begin{itemize}
\tightlist
\item
  A weakening or blocking of the upright force
\item
  The same energy turned inward, withheld, or self-directed
\item
  In negative cards: the worst has passed, or the difficulty is internal rather than external
\end{itemize}

Waite\textquotesingle s reversed meanings are specific --- they are not simply negations of the upright. The reversed Three of Cups is not the absence of joy but overindulgence. The reversed Ten of Swords is not the continuation of ruin but what comes after it. The reversed meanings in this document follow Waite\textquotesingle s original formulations in the Pictorial Key.

A practitioner coming from the Golden Dawn system, which uses elemental dignities rather than reversals, will find the WS reversal system cruder but faster. The dignities system is more precise for complex readings; reversals are more practical for single-card and three-card work.

\begin{center}\rule{0.5\linewidth}{0.5pt}\end{center}

\section{Notes on Combinations}\label{notes-on-combinations}

The Smith-Waite deck rewards reading scene to scene as much as card to card --- two adjacent scenes that share visual echoes (the water in the Six of Cups flowing into the landscape of the Seven of Cups; the blindfolded figure of the Two of Swords continuing into the imprisoned figure of the Eight) reinforce and modify each other in ways the abstracted GD meanings do not produce. Smith built a visual language that is partly sequential and partly associative. Attending to what scenes share --- color, figure posture, environmental condition, the direction figures face --- opens a layer of combination that operates independently of positional rules.

Beyond visual combination, the Sephirothic logic underlying the pip cards provides the same combinatorial structure as in the Golden Dawn system: a 6 next to a 4 describes Tiphareth in relation to Chesed --- harmony meeting foundation, the accomplished heart of the matter encountering what is stable and established. A practitioner with GD training can read the numeric combinations through their Sephirothic logic even when using Smith\textquotesingle s pictorial deck.
